% Generated draft: requires mathematical and source review.
\documentclass[11pt]{article}
\usepackage[margin=25mm]{geometry}
\usepackage{fontspec}
\setmainfont{Latin Modern Roman}
\usepackage{amsmath,amssymb,mathtools,unicode-math}
\setmathfont{Latin Modern Math}
\usepackage{enumitem,needspace,xurl,hyperref}
\setlength{\parindent}{0pt}
\setlength{\parskip}{5pt}
\setlength{\emergencystretch}{3em}
\newcommand{\N}{\mathbb{N}}
\newcommand{\Z}{\mathbb{Z}}
\newcommand{\Q}{\mathbb{Q}}
\newcommand{\R}{\mathbb{R}}
\newcommand{\C}{\mathbb{C}}
\newcommand{\F}{\mathbb{F}}
\newcommand{\A}{\mathbb{A}}
\newcommand{\PP}{\mathbb{P}}
\newcommand{\E}{\mathbb{E}}
\newcommand{\eps}{\varepsilon}
\newcommand{\dd}{\,\mathrm{d}}
\begin{document}
\section*{UnsolvedMath ID: 800917}
\textbf{Source problem code:} catalogue entry
\subsection{Definitions and mathematical statement}
Let \(d\) denote the dimension of an ambient algebraic variety \(X\), let \(E\) be a sheaf represented by a proposed moduli space \(M\), and write \(\operatorname{Ext}^{i}_{X}(E,E)\) for its global self-extension groups. The source describes \(\operatorname{Ext}^{1}_{X}(E,E)\) as governing deformations and \(\operatorname{Ext}^{2}_{X}(E,E)\) as governing obstructions. It gives, explicitly only as a rough requirement for constructing a virtual fundamental class, the vanishing condition
\[
\operatorname{Ext}^{i}_{X}(E,E)=0\qquad\text{for every integer }i>2.
\]
This explanatory condition is not supplied as a precise necessary-and-sufficient criterion or as an additional hypothesis on the requested higher-dimensional theory. The source says that on threefolds suitable vanishing can often be obtained using Serre duality and stability, without specifying the hypotheses or the version of the deformation complex involved.

\textbf{Problem (Pandharipande, 2010, Q5).} Define Donaldson--Thomas invariants in dimensions \(d>3\).

This is a construction problem: develop a higher-dimensional theory of counting sheaves of Donaldson--Thomas type. The supplied formulation does not specify the ground field, geometric hypotheses on \(X\), class or stability condition of the sheaves, numerical data, properness conditions, obstruction-theoretic construction, or values in which the proposed invariants should lie. It also does not specify whether the goal concerns every dimension greater than three or particular higher-dimensional settings. These choices remain unresolved; no assertion of existence for every variety or every sheaf moduli space is intended by this draft.

As motivation, the source cites higher-dimensional enumerative computations in Gromov--Witten theory under the labels [klp,pz] and says that many aspects of Joyce's counting theory apply in higher dimensions [J]. Its introductory assertion that Donaldson--Thomas invariants are defined only in dimension three is retained as historical source context, not as an independently verified present-day claim.
\subsection{Short English statement}
Pandharipande asks for a definition of Donaldson--Thomas invariants in dimensions greater than three. The source motivates the question through deformation and obstruction groups of sheaves, but leaves the geometric setting and the required properties of the invariants unspecified. The supplied registry labels the problem open; this draft does not independently verify that status.
\subsection{Sources}
\begin{itemize}
\item Local UnsolvedMath registry, permanent ID \(800917\), exact source problem code \texttt{catalogue entry}; supplied title: \emph{AGNES Open Problems in Algebraic Geometry --- Q5}. The record attributes the problem to Pandharipande, with proposed year 2010, and labels it open. The supplied record was updated on 2026-10-16; this is registry metadata, not independent verification of the mathematical status.
\item Pandharipande, \emph{Open Problems in Enumerative Algebraic Geometry (AGNES)} (2010), source item Q5 (\#4), \url{https://arxiv.org/abs/1004.3259}. The supplied provenance records access on 2026-10-16 and extraction method \texttt{arxiv\_latex}. This draft uses the supplied extraction and does not claim to have visited or checked the linked source.
\item Source reference labels [klp,pz]: cited for parallel examples of higher-dimensional enumerative computations in Gromov--Witten theory. No bibliographic entries or corresponding links are supplied.
\item Source reference label [J]: cited for the assertion that many aspects of Joyce's counting theory are valid in higher dimensions. No bibliographic entry or corresponding link is supplied.
\end{itemize}
\end{document}
